\section{Methodology}
\label{sec:methodology}

Building on our observation that accuracy degradation under distribution shift may be uniform across models, we design a methodology to directly measure ranking stability between ImageNet and ImageNet-V2.

\subsection{Overview}

Our approach has three components: (1) data collection from standardized leaderboards, (2) computation of ranking correlation statistics, and (3) bootstrap inference for uncertainty quantification.

\subsection{Data Collection}

\textbf{Source.} We collect model accuracy data from Papers With Code leaderboards for ImageNet and ImageNet-V2 \cite{paperswithcode}. Papers With Code provides community-curated, standardized results using consistent evaluation protocols.

\textbf{Inclusion Criteria.} We include all models that have reported results on both ImageNet (ILSVRC 2012 validation set) and ImageNet-V2 (MatchedFrequency variant). This yielded 96 models spanning publication years 2015--2024 and architecture families including ResNet, ViT, ConvNeXt, EfficientNet, and others.

\textbf{Rationale.} Using leaderboard data ensures comparable evaluation protocols across models. While selection bias may exist (models with poor V2 results may be underreported), this bias likely understates ranking instability---our finding of high stability is thus conservative.

\subsection{Ranking Computation}

For each benchmark, we rank models by top-1 accuracy in descending order (higher accuracy = better rank). Ties are handled using standard ranking conventions.

\subsection{Kendall-$\tau$ Correlation}

We compute Kendall-$\tau$ (tau-b) between the two ranking vectors. Kendall-$\tau$ measures the proportion of concordant vs. discordant pairs.

\textbf{Rationale.} Kendall-$\tau$ is preferred for ordinal rankings because it (1) is robust to ties, (2) has a natural interpretation as proportion of correctly ordered pairs, and (3) is less sensitive to outliers than Pearson correlation.

\textbf{Threshold.} We pre-registered $\tau < 0.90$ as ``significant ranking shift'' based on prior work suggesting $\tau \geq 0.95$ indicates essentially identical rankings.

\subsection{Bootstrap Confidence Intervals}

We estimate 95\% confidence intervals via bootstrap resampling (10,000 iterations) using the percentile method \cite{efron1993bootstrap}.

\subsection{Hypothesis Testing}

We evaluate our pre-registered hypothesis:

\begin{itemize}
    \item \textbf{H1:} $\tau < 0.90$ (rankings shift significantly)
    \item \textbf{H0:} $\tau \geq 0.90$ (rankings are preserved)
\end{itemize}

If the 95\% CI upper bound falls below 0.90, we reject H0 in favor of H1. If $\tau \geq 0.90$ and the CI is above 0.90, we fail to reject H0 and conclude rankings are preserved.
